%% kapa: exam=1 round=35 year=2024 period=1 form=A subject=real_estate subject_ko=부동산학원론 seq=17 number=97
%% source: Q-Net 감정평가사 기출문제 (cst00309) artlSeq=5213229 file=raw/1008/5213229/1교시 문제지 원본.pdf
%% conversion: kordoc 4.24.51 -> tools/gen1.py (auto); review=proofread-sonnet; flags=none
\begin{question}{97}
감정평가의 지역분석에 관한 내용으로 옳은 것은?
\begin{choices}
  \item 인근지역이란 감정평가의 대상이 된 부동산이 속한 지역으로서 부동산의 이용이 동질적이고 가치형성요인 중 지역요인을 공유하는 지역을 말한다.
  \item 유사지역이란 대상부동산이 속한 지역으로서 인근지역과 유사한 특성을 갖는 지역을 말한다.
  \item 동일수급권이란 대상부동산과 수요ㆍ공급 관계가 성립하고 가치 형성에 서로 영향을 미치지 않는 관계에 있는 다른 부동산이 존재하는 권역을 말한다.
  \item 지역분석은 대상지역 내 토지의 최유효이용 및 대상부동산의 가격을 판정하는 것이다.
  \item 지역분석은 개별분석 이후에 실시하는 것이 일반적이다.
\end{choices}
\end{question}
\answer{1}
